\chapter{Things and Arrows}

% ═══════════════════════════════════════════════════════════════════════════
\section{What Category Theory Is About}
\label{sec:what-ct-is-about}
% ═══════════════════════════════════════════════════════════════════════════

\CoreLabel
\subsection{Core: The Central Idea}

Category theory is the study of \term{things} and \term{arrows between things}.

That's it. That's the whole subject.

Of course, ``things'' and ``arrows'' are deliberately vague words. The power
of category theory comes precisely from that vagueness. A thing can be a
recipe, a city, a sentence, a number, a proof, or a process. An arrow can be
a road, a translation, a logical implication, a comparison, or a
transformation of any kind.

What category theory cares about is not \emph{what things are made of
inside}, but \emph{how they connect to other things}. Two things that connect
to everything else in exactly the same way are, for the purposes of category
theory, \emph{the same thing}---even if they look completely different
up close.

This is a radical idea, and we will return to it many times.

% ─── Expansion 1 ────────────────────────────────────────────────────────────
\ExpOneLabel
\subsection{Expansion 1: Three Everyday Pictures}

Let's build the intuition with three examples that have nothing to do with
mathematics.

\begin{example}{The Transit Map}
A city has train stations. Between stations there are train lines. If you want
to get from Station~A to Station~C, and there is a line from A to B and a line
from B to C, you can travel A~$\to$~B~$\to$~C.

\medskip
\begin{center}
\begin{tikzcd}[column sep=2.5em]
  A \arrow[r, "\text{line 1}"] & B \arrow[r, "\text{line 2}"] & C
\end{tikzcd}
\end{center}
\medskip

The \emph{things} are the stations. The \emph{arrows} are the train lines.
Notice: we don't care what is inside a station---its architecture, its size,
whether it has a coffee shop. We only care how it connects.
\end{example}

\begin{dialog}
When I see a diagram like the one above, I ask myself three questions:
\begin{enumerate}
  \item \emph{What are the dots?} Here: train stations.
  \item \emph{What are the arrows?} Here: train lines you can ride.
  \item \emph{Can I follow a path?} Here: yes, A to B to C.
\end{enumerate}
These three questions apply to \emph{every} diagram in this book. Make them
into a habit.
\end{dialog}

\begin{example}{The Recipe Steps}
A recipe is a sequence of steps. ``Chop onions'' must happen before ``sauté
onions.'' ``Sauté onions'' must happen before ``add to pot.''

\medskip
\begin{center}
\begin{tikzcd}[column sep=2.5em]
  \text{chop} \arrow[r, "\text{before}"] &
  \text{sauté} \arrow[r, "\text{before}"] &
  \text{add to pot}
\end{tikzcd}
\end{center}
\medskip

The \emph{things} are the steps. The \emph{arrows} mean ``must happen before.''
If chopping must happen before sautéing, and sautéing before adding, then
chopping must happen before adding---even though there is no single arrow
drawn directly from ``chop'' to ``add to pot.''

This indirect conclusion is not magic. It follows from the arrows we have.
In category theory we will give this a formal name: \term{composition}.
\end{example}

\begin{example}{The Approval Chain}
A company requires that budget requests be approved by a manager before going
to the director, and by the director before going to the board.

\medskip
\begin{center}
\begin{tikzcd}[column sep=2.5em]
  \text{request} \arrow[r, "\text{manager}"] &
  \text{approved} \arrow[r, "\text{director}"] &
  \text{final approval}
\end{tikzcd}
\end{center}
\medskip

Again: \emph{things} (states of the request), \emph{arrows} (approvals that
move the request forward). The pattern is identical to the transit map and the
recipe. This is the first sign that something deeper is going on: very
different situations share the same shape.

Category theory is, among other things, the science of recognizing when two
situations have the same shape.
\end{example}

% ─── Expansion 2 ────────────────────────────────────────────────────────────
\ExpTwoLabel
\subsection{Expansion 2: Inside vs.\ Outside}

There is a habit of mind that category theory requires, and it is worth
practicing from the very beginning.

\textbf{Stop asking what things are made of.}

When you look at a train station on a map, you do not ask how many
platforms it has or which architect designed it. You ask: which lines does it
connect to? Category theory extends this discipline to \emph{everything}. When
we work with a thing in category theory, we are forbidden from looking
``inside'' it. We can only look at what arrows arrive at it and what arrows
leave it.

This sounds limiting. In practice it is liberating.

Here is why. Suppose you have two completely different-looking objects---a
subway station and an airport terminal, say---but every route that goes through
one also goes through the other, and vice versa. For the purposes of your
journey, they are \emph{interchangeable}. Category theory gives you a way to
say that precisely and work with it rigorously.

\begin{readernote}
Try this: pick any two things in your life that are functionally
interchangeable but look different inside. A paperback and an e-book of the
same novel. An email and a printed letter carrying the same message. A
physical key and a key card opening the same door.

In each case: the \emph{inside} differs, but the \emph{connections} (what
can be done with the object, where it fits) are the same. In category theory
we say these are \term{isomorphic}---a word we will define carefully later.
\end{readernote}

The discipline of ``look only at connections, not at insides'' is called
working \term{externally} or working \term{up to isomorphism}. We will build
to a precise definition. For now, just start training the habit: when you
look at a thing, ask not \emph{what is it?} but \emph{what does it connect
to?}

% ─── Expansion 3 ────────────────────────────────────────────────────────────
\ExpThreeLabel
\subsection{Expansion 3: Historical and Philosophical Note}

Category theory was introduced in 1945 by Samuel Eilenberg and Saunders
Mac~Lane in a paper titled ``General Theory of Natural Equivalences.'' Their
immediate goal was to make precise a certain kind of structural similarity they
kept encountering in algebraic topology. But the language they invented turned
out to be far more general.

The philosopher's way to describe what category theory does is this: it
replaces the question \emph{``what is X?''} with the question \emph{``what
role does X play?''} This is sometimes called \term{structuralism} in the
philosophy of mathematics---the view that mathematical objects have no
intrinsic identity, only relational identity.

From the structuralist perspective, there is no single ``correct'' natural
number~2. There are many different formal systems in which something plays the
role of~2. What makes them all \emph{the} number 2 is that they each stand in
the same relationships to the things around them.

This book does not require you to adopt a philosophical position. But it is
worth knowing that category theory carries a viewpoint with it, and that
viewpoint has consequences for how we write definitions, how we state theorems,
and what we count as ``solving'' a problem.


% ═══════════════════════════════════════════════════════════════════════════
\section{Objects and Morphisms}
\label{sec:objects-morphisms}
% ═══════════════════════════════════════════════════════════════════════════

\CoreLabel
\subsection{Core: Naming the Pieces}

The \emph{things} in category theory have an official name: \term{objects}.

The \emph{arrows} have an official name: \term{morphisms}.

Every morphism goes \emph{from} one object \emph{to} another. We write
\[
  f : A \to B
\]
to mean ``there is a morphism named~$f$ going from object~$A$ to object~$B$.''
We call $A$ the \term{domain} of~$f$ and $B$ the \term{codomain} of~$f$.

That's the whole vocabulary for now: objects, morphisms, domain, codomain.
Everything else in category theory is built from these four words.

% ─── Expansion 1 ────────────────────────────────────────────────────────────
\ExpOneLabel
\subsection{Expansion 1: Reading Morphism Notation}

Let's slow down and practice reading the notation $f : A \to B$, because it
will appear on almost every page of this book.

\begin{dialog}
I see the expression $f : A \to B$. Here is what I say to myself:

\begin{itemize}
  \item ``$f$ is the name of a morphism.''
  \item ``$A$ is where $f$ starts---its domain.''
  \item ``$B$ is where $f$ ends---its codomain.''
  \item ``The arrow $\to$ tells me this is a one-directional connection.''
  \item ``I do \emph{not} assume I can go from $B$ back to $A$ just because
        I can go from $A$ to $B$. Arrows have directions.''
\end{itemize}

If someone asks me to draw this, I draw a dot labeled $A$, a dot labeled $B$,
and an arrow from $A$ to $B$ labeled $f$:

\medskip
\begin{center}
\begin{tikzcd}
  A \arrow[r, "f"] & B
\end{tikzcd}
\end{center}
\medskip

This picture and the expression $f : A \to B$ carry exactly the same
information. I can translate between them instantly.
\end{dialog}

Let's look at some concrete examples with this vocabulary in place.

\begin{example}{Morphisms as Comparisons}
Suppose our objects are temperature readings: $32°$, $68°$, $98°$, $212°$.
A morphism $f : 32° \to 212°$ could mean ``is colder than.'' A morphism
$g : 68° \to 98°$ could mean ``is colder than.'' Notice that both $f$ and $g$
are instances of the same \emph{kind} of arrow (``is colder than''), but they
are different morphisms because they connect different objects.

The label of a morphism and its domain and codomain together make it what it
is. Two morphisms are the same only if they have the same name \emph{and} the
same domain \emph{and} the same codomain.
\end{example}

\begin{example}{Morphisms as Translations}
Suppose our objects are the phrases ``hello,'' ``hola,'' ``bonjour,'' ``ciao.''
A morphism $t : \text{hello} \to \text{hola}$ could mean ``is a Spanish
translation of.'' A morphism $u : \text{hello} \to \text{bonjour}$ could mean
``is a French translation of.''

Both start at the same object (``hello'') but end at different objects. They
are different morphisms. The object ``hello'' has \emph{two outgoing
morphisms} in this setup.
\end{example}

\begin{readernote}
When you encounter a new example, before doing anything else, identify:
\begin{enumerate}
  \item What are the \emph{objects}? (the dots in the diagram)
  \item What are the \emph{morphisms}? (the arrows)
  \item For each morphism: what is its \emph{domain}? what is its
        \emph{codomain}?
\end{enumerate}
Write these down, in words, every time. This routine takes thirty seconds and
prevents ninety percent of the confusion that comes from working with
diagrams too quickly.
\end{readernote}

% ─── Expansion 2 ────────────────────────────────────────────────────────────
\ExpTwoLabel
\subsection{Expansion 2: What Morphisms Are Not}

Several natural assumptions about morphisms turn out to be wrong in general.
Let's eliminate them now before they become habits.

\textbf{Morphisms are not always functions.}
In many examples, morphisms will be functions from one collection to another.
But the word ``morphism'' does not mean ``function.'' A morphism is anything
that fits the pattern: it has a domain, a codomain, and it goes from one to
the other in a way consistent with the rules of its category. We will see
morphisms that are proofs, comparisons, paths, transformations, and more.

\textbf{Morphisms can have the same domain and codomain without being equal.}
If $f : A \to B$ and $g : A \to B$, that does \emph{not} mean $f = g$. There
can be many different morphisms between the same two objects. In the
translation example above, if we added a second Spanish translation (say, a
formal one and an informal one), we'd have two distinct morphisms from
``hello'' to ``hola.''

\textbf{A morphism $f : A \to B$ does not imply a morphism $B \to A$.}
Going from $A$ to $B$ is not the same as going from $B$ to $A$, and the
existence of one does not guarantee the existence of the other. In the recipe
example, ``chop must happen before sauté'' does not imply ``sauté must happen
before chop''---in fact, that would be a contradiction.

\textbf{There may be no morphism between two objects.}
In any given category, some pairs of objects may simply not be connected.
That's fine. Not every city has a direct train to every other city.

% ─── Expansion 3 ────────────────────────────────────────────────────────────
\ExpThreeLabel
\subsection{Expansion 3: Notation and Conventions}

\deepexpandsection{Alternative Notation}

In some literature you will see morphisms written differently:
\begin{itemize}
  \item $f \in \mathrm{Hom}(A, B)$ --- reads ``$f$ is a member of the
        \term{hom-set} from $A$ to $B$.'' The hom-set $\mathrm{Hom}(A,B)$
        is the collection of all morphisms with domain $A$ and codomain $B$.
  \item $A \xrightarrow{f} B$ --- an inline arrow notation, used when we
        want to show a morphism inside a sentence without breaking into
        display math.
  \item $\mathrm{Hom}_{\mathcal{C}}(A, B)$ --- the subscript $\mathcal{C}$
        specifies \emph{which category} we mean, since the same objects can
        appear in multiple categories with different morphism sets.
\end{itemize}

In this book we default to $f : A \to B$ and introduce the hom-set notation
when we need to talk about the whole collection of morphisms between two
objects at once.

\deepexpandsection{The Category of This Section's Examples}

The temperature example above is almost---but not quite---a category. We have
objects ($32°$, $68°$, $98°$, $212°$) and we have morphisms (the ``colder
than'' relations). Two things are missing: we need \emph{identity morphisms}
and we need \emph{composition}. We will add both in the next two sections and
then the example will be a full category.


% ═══════════════════════════════════════════════════════════════════════════
\section{Composition}
\label{sec:composition}
% ═══════════════════════════════════════════════════════════════════════════

\CoreLabel
\subsection{Core: Chaining Arrows}

If there is a morphism from $A$ to $B$, and a morphism from $B$ to $C$, then
there must also be a morphism from $A$ to $C$.

This is the rule of \term{composition}. It says: whenever two arrows line up
end-to-end, you get a third arrow for free---the one that follows the same
path from start to finish.

\medskip
\begin{center}
\begin{tikzcd}[column sep=3em, row sep=2em]
  A \arrow[r, "f"] \arrow[rr, bend right=25, "g \circ f"'] & B \arrow[r, "g"] & C
\end{tikzcd}
\end{center}
\medskip

We write the composed morphism as $g \circ f$, read ``$g$ after $f$'' or
``$g$ composed with $f$.'' It goes from $A$ to $C$.

Notice the order: $g \circ f$ is read right to left. First $f$, then $g$.
This will feel backwards at first. We'll practice until it feels natural.

% ─── Expansion 1 ────────────────────────────────────────────────────────────
\ExpOneLabel
\subsection{Expansion 1: Practicing Composition}

\begin{dialog}
I see:
\[
  f : A \to B \qquad g : B \to C
\]
I ask: ``Do these line up?'' The codomain of $f$ is $B$. The domain of $g$
is $B$. Yes, they match. So composition is possible.

I write the composed morphism as $g \circ f : A \to C$.

I read this as: ``first do $f$ (go from $A$ to $B$), then do $g$ (go from
$B$ to $C$). The overall result goes from $A$ to $C$.''

I always check: domain of composed = domain of first, codomain of composed
= codomain of last.
\[
  (g \circ f) : A \to C
  \qquad\text{domain = }A,\quad\text{codomain = }C.
\]
\end{dialog}

\begin{example}{Composition in the Recipe}
Recall our recipe morphisms:
\[
  f : \text{chop} \to \text{sauté}, \qquad g : \text{sauté} \to \text{add to pot}
\]
Composition gives:
\[
  g \circ f : \text{chop} \to \text{add to pot}
\]
This says: ``chopping must happen before adding to pot.'' We knew that
intuitively. Composition is the formal way to say it.
\end{example}

\begin{example}{A Chain of Three}
Suppose we have:
\[
  f : A \to B, \qquad g : B \to C, \qquad h : C \to D
\]
We can compose step by step:
\[
  g \circ f : A \to C, \qquad h \circ (g \circ f) : A \to D
\]
Or we can compose from the other end:
\[
  h \circ g : B \to D, \qquad (h \circ g) \circ f : A \to D
\]

\medskip
\begin{center}
\begin{tikzcd}[column sep=2.5em]
  A \arrow[r, "f"] \arrow[rrr, bend right=35, "h \circ g \circ f"']
  & B \arrow[r, "g"] \arrow[rr, bend left=30, "h \circ g"]
  & C \arrow[r, "h"]
  & D
\end{tikzcd}
\end{center}
\medskip

Both routes give the same composed morphism from $A$ to $D$. This is not a
coincidence---it is a requirement called \term{associativity}, which we will
examine closely in Expansion~2 below.
\end{example}

\begin{readernote}
Every time you see two arrows that share a middle object, ask: ``What does
their composition tell me?'' Draw the composed arrow (perhaps with a dashed
line or a curved arrow) and label it. This keeps your diagram annotated and
your thinking explicit.
\end{readernote}

% ─── Expansion 2 ────────────────────────────────────────────────────────────
\ExpTwoLabel
\subsection{Expansion 2: Associativity}

Associativity means that when you compose three morphisms in a chain, it
doesn't matter which composition you do first---you always get the same result.

\[
  h \circ (g \circ f) = (h \circ g) \circ f
\]

Why does this matter? Because it means we can write $h \circ g \circ f$
without parentheses and there is no ambiguity. The chain of three arrows has a
unique composite, regardless of how we group the steps.

\begin{example}{Associativity in Transit}
Suppose:
\begin{itemize}
  \item $f$: travel from Airport to Downtown
  \item $g$: travel from Downtown to Stadium
  \item $h$: travel from Stadium to Hotel
\end{itemize}
The trip ``Airport $\to$ Downtown $\to$ Stadium $\to$ Hotel'' is the same
trip whether you think of it as:
\begin{itemize}
  \item (Airport $\to$ Downtown-Stadium) then Stadium $\to$ Hotel, or
  \item Airport $\to$ Downtown then (Downtown $\to$ Stadium-Hotel).
\end{itemize}
Grouping the \emph{planning} differently does not change the \emph{journey}.
Associativity just says this obvious fact is always true for morphisms.
\end{example}

Associativity has a deeper consequence that is easy to miss: it means that any
\emph{finite chain} of composable arrows has a unique composite. We don't need
to specify how to parenthesize it. The composite of $f_1, f_2, \ldots, f_n$
(when they chain up in order) is well-defined.

% ─── Expansion 3 ────────────────────────────────────────────────────────────
\ExpThreeLabel
\subsection{Expansion 3: Composition as the Primary Structure}

\deepexpandsection{Why Composition Comes Before Identity}

Many presentations introduce identity morphisms before composition. We
reversed the order deliberately.

Composition is the \emph{operative} structure of a category---it is what
makes the arrows into more than just a list of isolated connections. Identity
morphisms, as we will see in the next section, are a \emph{constraint} on
composition: they are the neutral elements that leave composition unchanged.
Logically, you must understand the operation before you can understand its
neutral element.

\deepexpandsection{Partial vs.\ Total Composition}

In a category, composition is defined \emph{only} for pairs of morphisms where
the codomain of the first equals the domain of the second. It is a
\emph{partial} operation on morphisms, not a total one. This is different from
addition of numbers, which is defined for every pair of numbers.

Formally: given morphisms $f : A \to B$ and $g : C \to D$, the composite
$g \circ f$ is defined \emph{if and only if} $B = C$ (the codomain of $f$
equals the domain of $g$).

This is why the domain and codomain of each morphism are so important to track.
They are the ``type-checking'' mechanism that tells you which compositions are
legal.


% ═══════════════════════════════════════════════════════════════════════════
\section{Identity Morphisms}
\label{sec:identity}
% ═══════════════════════════════════════════════════════════════════════════

\CoreLabel
\subsection{Core: Every Object Has a ``Stay Put'' Arrow}

Every object has a special morphism that goes from \emph{itself to itself} and
does nothing. We call it the \term{identity morphism} on that object and write
$\mathrm{id}_A : A \to A$.

``Does nothing'' has a precise meaning: composing any morphism with the
identity gives back the original morphism unchanged.

\begin{center}
\begin{tikzcd}[column sep=3em]
  A \arrow[r, "f"] \arrow[loop left, "\mathrm{id}_A"]
  & B \arrow[loop right, "\mathrm{id}_B"]
\end{tikzcd}
\end{center}

For any morphism $f : A \to B$:
\[
  f \circ \mathrm{id}_A = f \qquad \text{and} \qquad \mathrm{id}_B \circ f = f
\]

Doing ``nothing'' and then $f$ gives $f$. Doing $f$ and then ``nothing'' gives
$f$. The identity is the neutral element of composition.

% ─── Expansion 1 ────────────────────────────────────────────────────────────
\ExpOneLabel
\subsection{Expansion 1: Seeing Identity in Everyday Situations}

Identity morphisms can feel like a technicality. They are not. Let's see why
they appear naturally.

\begin{example}{The ``Do Nothing'' Step in a Recipe}
Suppose you write a recipe that can optionally rest the dough for an hour. On
a day when you're in a hurry, you do \emph{nothing} during that step---you
skip it, and the dough is the same dough it was before.

This ``do nothing'' is an identity morphism: it takes the dough as input
and returns the dough unchanged. Every state that can appear in a process has
this ``pass through unchanged'' option.
\end{example}

\begin{example}{The Zero-Distance Route}
In the transit example, ``stay at Station A'' is a perfectly valid route: you
start at A and end at A. It composes correctly with other routes:
\[
  (\text{stay at A}) \text{ then } (\text{A to B}) = (\text{A to B})
\]
\[
  (\text{A to B}) \text{ then } (\text{stay at B}) = (\text{A to B})
\]
The ``stay put'' route is the identity morphism on Station A.
\end{example}

\begin{dialog}
When I see a loop arrow on an object---an arrow that starts and ends at the
same dot---my first question is: ``Is this the identity, or is it a
non-trivial self-loop?''

The identity is special because \emph{composing with it does nothing}. A
non-trivial self-loop (there can be others) might actually change something.

I can always tell them apart by testing: does composing with this arrow change
any other morphism? If not, it's the identity. If yes, it's something else.
\end{dialog}

% ─── Expansion 2 ────────────────────────────────────────────────────────────
\ExpTwoLabel
\subsection{Expansion 2: Uniqueness of Identity}

Each object has \emph{exactly one} identity morphism. This is not an
assumption---it's a theorem. Here's the argument:

Suppose $\mathrm{id}_A$ and $\mathrm{id}_A'$ are both identity morphisms on
$A$ (i.e., both satisfy the neutrality law). Then:
\[
  \mathrm{id}_A
  = \mathrm{id}_A \circ \mathrm{id}_A'
  = \mathrm{id}_A'
\]
The first equality uses the fact that $\mathrm{id}_A'$ is a right identity.
The second uses the fact that $\mathrm{id}_A$ is a left identity. Both are
the same morphism.

This is the first example in this book of a categorical proof: short,
self-referential, and surprisingly powerful. We will see many more like it.

\begin{readernote}
Notice the structure of the proof: we used the identity \emph{laws} (equations
that must hold) rather than looking inside any object. That's the categorical
style. If a proof ever asks you to look inside an object, something has gone
wrong.
\end{readernote}

% ─── Expansion 3 ────────────────────────────────────────────────────────────
\ExpThreeLabel
\subsection{Expansion 3: Identity as a Function, Formally}

\deepexpandsection{Identity as an Assignment}

The identity is not just ``a morphism per object''---it is an assignment:
given any object $A$, we assign the morphism $\mathrm{id}_A : A \to A$.
This assignment is part of the \emph{data} of a category. When we define a
category formally (in the next chapter), we will include this assignment in
the definition.

\deepexpandsection{Connection to Monoids}

For readers who encounter the word later: a \term{monoid} is a set with an
associative binary operation and an identity element. The morphisms from a
single object $A$ to itself, together with composition, form a monoid. The
identity morphism is the identity element of that monoid. This connection
between single-object categories and monoids is one of the first of many
``hidden structures'' that category theory uncovers---but we will build to it
properly rather than assuming it here.


% ═══════════════════════════════════════════════════════════════════════════
\section{Commutative Diagrams}
\label{sec:commutative-diagrams}
% ═══════════════════════════════════════════════════════════════════════════

\CoreLabel
\subsection{Core: When All Paths Agree}

A \term{commutative diagram} is a diagram of objects and morphisms in which
every path between the same two endpoints gives the same composite morphism.

The simplest example is a triangle:

\medskip
\begin{center}
\begin{tikzcd}[column sep=3em, row sep=2.5em]
  A \arrow[r, "f"] \arrow[dr, "h"'] & B \arrow[d, "g"] \\
  & C
\end{tikzcd}
\end{center}
\medskip

This diagram \emph{commutes} if $g \circ f = h$.

There are two paths from $A$ to $C$: the path $A \to B \to C$ (composing $f$
then $g$) and the direct path $A \to C$ (the morphism $h$). Commutativity
says these give the same result.

% ─── Expansion 1 ────────────────────────────────────────────────────────────
\ExpOneLabel
\subsection{Expansion 1: Reading a Commutative Diagram}

\begin{dialog}
I see a diagram with multiple paths between two objects. My procedure:

\begin{enumerate}
  \item \textbf{Identify} all the paths from each object to each other object.
        I trace every possible route, like reading a map.
  \item \textbf{Write} the composite morphism for each path, using the
        $\circ$ notation.
  \item \textbf{Ask}: are these composites equal? If yes, the diagram
        (or at least those paths) commutes.
  \item \textbf{Note} that a square diagram has two paths from top-left to
        bottom-right, a pentagon has multiple paths from each vertex, and so on.
\end{enumerate}
\end{dialog}

\begin{example}{A Commutative Square}
The most common diagram in category theory is the commutative square:

\medskip
\begin{center}
\begin{tikzcd}[column sep=3.5em, row sep=2.5em]
  A \arrow[r, "f"] \arrow[d, "p"'] & B \arrow[d, "q"] \\
  C \arrow[r, "g"'] & D
\end{tikzcd}
\end{center}
\medskip

This commutes when $q \circ f = g \circ p$.

In words: going right then down gives the same result as going down then right.

\medskip

Commutative square in everyday terms: suppose $A$ is a sentence in English,
$B$ is a formal logical statement, $C$ is the same sentence translated into
French, and $D$ is the same formal logical statement. Then:
\begin{itemize}
  \item $f$: formalize the English sentence.
  \item $p$: translate the English sentence to French.
  \item $q$: formalize the French sentence.
  \item $g$: translate the English formalization to French formalization
        (or: verify they match).
\end{itemize}
The square commutes if ``translate then formalize'' gives the same result as
``formalize then translate.'' Whether this holds depends on the specific
language and logic in question---it is a genuine mathematical question, not
a foregone conclusion.
\end{example}

\begin{readernote}
Whenever you encounter a diagram in this book (or anywhere in mathematics),
write next to it: ``This commutes because \_\_\_\_\_\_'' and fill in the blank.
If you can't fill it in, you don't yet understand the claim the diagram is
making. Diagrams are \emph{assertions}---they say something is equal. Read
them as equations, not just pictures.
\end{readernote}

% ─── Expansion 2 ────────────────────────────────────────────────────────────
\ExpTwoLabel
\subsection{Expansion 2: Diagrams as Equations}

Every commutative diagram is shorthand for a collection of equations between
composed morphisms. The diagram is often cleaner to read than the equations,
which is why we use it. But the two are entirely equivalent.

A triangle diagram asserts one equation. A square diagram asserts one
equation (two paths from top-left to bottom-right). A more complex diagram
with five or six objects might assert many equations simultaneously, one for
each pair of objects with multiple paths between them.

This is the real power of diagrammatic notation: a single picture can
communicate a system of equations in a way that is easier to parse than
writing all the equations in a list. As you read more category theory, you
will find yourself preferring the diagram and translating to equations only
when you need to compute.

There is also a converse skill: when you have a collection of equations
relating composed morphisms, you can often \emph{draw} them as a commutative
diagram. Practicing this translation in both directions deepens understanding.

% ─── Expansion 3 ────────────────────────────────────────────────────────────
\ExpThreeLabel
\subsection{Expansion 3: Diagrams in Formal Proofs}

\deepexpandsection{The Pasting Lemma}

One of the most useful facts about commutative diagrams is the \term{pasting
lemma}: if you have two commutative squares sharing an edge, the outer
rectangle also commutes.

\medskip
\begin{center}
\begin{tikzcd}[column sep=3em, row sep=2.5em]
  A \arrow[r, "f"] \arrow[d, "p"']
  & B \arrow[r, "g"] \arrow[d, "q"]
  & C \arrow[d, "r"] \\
  D \arrow[r, "s"']
  & E \arrow[r, "t"']
  & F
\end{tikzcd}
\end{center}
\medskip

If the left square commutes ($q \circ f = s \circ p$) and the right square
commutes ($r \circ g = t \circ q$), then the outer rectangle commutes
($r \circ g \circ f = t \circ s \circ p$).

This lemma is used constantly in categorical proofs. Once you know it, you
can build large commutative diagrams from smaller confirmed pieces, like
tiling a floor.

\deepexpandsection{Diagrams Without Commutativity}

Not every diagram in category theory is required to commute. Sometimes we draw
a diagram just to display objects and morphisms without asserting any
equations. The convention in this book: if a diagram is stated to commute
(either explicitly or by context), it does. Otherwise, assume it is merely
depicting the structure.

When in doubt, check whether the author says ``the following diagram commutes''
or just ``consider the following diagram.'' This distinction matters.


% ═══════════════════════════════════════════════════════════════════════════
\section{Chapter Summary}
\label{sec:ch1-summary}
% ═══════════════════════════════════════════════════════════════════════════

\begin{summarybox}
\textbf{Objects} are the things in a category. They have no specified internal
structure---we work with them only through their connections.

\textbf{Morphisms} ($f : A \to B$) are the arrows between objects. Each
morphism has a domain and a codomain.

\textbf{Composition}: if $f : A \to B$ and $g : B \to C$, then $g \circ f :
A \to C$ is their composite. Composition is associative: $(h \circ g) \circ f
= h \circ (g \circ f)$.

\textbf{Identity}: every object $A$ has a morphism $\mathrm{id}_A : A \to A$
that is neutral under composition: $f \circ \mathrm{id}_A = f$ and
$\mathrm{id}_B \circ f = f$ for any $f : A \to B$.

\textbf{Commutative diagrams} assert that all paths between the same two
objects compose to the same morphism. They are equations in disguise.

\medskip
\textbf{The reader's three questions}, to be asked of every new example:
\begin{enumerate}
  \item What are the objects?
  \item What are the morphisms (domain, codomain, name)?
  \item Which diagrams commute, and what equations do they assert?
\end{enumerate}
\end{summarybox}

\bigskip

In the next chapter, we assemble all these pieces---objects, morphisms,
composition, and identity---into a single package: a \term{category}.
